\subsection{Symplecticity, order and phase accuracy}
\label{sec:core-integrator-analysis}
For a differentiable local solution of~\eqref{eq:central4}, set
$C_\Delta=J\mathcal S_\Delta''(m)/2$. Differentiating the implicit map gives
\begin{equation}
 D\Psi_\Delta=(I-C_\Delta)^{-1}(I+C_\Delta).
 \label{eq:core-step-tangent}
\end{equation}
Since $\mathcal S_\Delta''$ is symmetric,
$C_\Delta^{\mathsf T}J+JC_\Delta=0$, so this Cayley transform satisfies
$(D\Psi_\Delta)^{\mathsf T}JD\Psi_\Delta=J$.
The odd dependence $\mathcal S_{-\Delta}=-\mathcal S_\Delta$ also gives
$\Psi_{-\Delta}=\Psi_\Delta^{-1}$. The local defect $O(\Delta^5)$ in
\eqref{eq:core-exact-displacement} implies fourth-order convergence on
fixed intervals with bounded derivatives and locally unique implicit
steps. These classical properties \citep[Theorems~17--18 and 23--24]{feng}
hold for the complete fitted energy; the computation below measures the
additional error from finite Newton solves.

The linear oscillator makes the phase improvement explicit. For
$h=\nu(X^2+Y^2)/2$, put $a=\nu\Delta+(\nu\Delta)^3/12$.
The step is a rotation with eigenvalues
$(1\pm\mathrm ia/2)/(1\mp\mathrm ia/2)$ and phase magnitude
\begin{equation}
 \theta_\Delta=2\arctan(a/2)
 =\nu\Delta-\frac{(\nu\Delta)^5}{120}+O(\Delta^7).
 \label{eq:core-linear-phase}
\end{equation}
Its modulus is one for every real step. Implicit midpoint instead has
$\theta_\Delta=\nu\Delta-(\nu\Delta)^3/12+O(\Delta^5)$.
Both preserve the quadratic energy exactly while accumulating different
phase errors. For a nonlinear energy, symplecticity preserves canonical
structure and does not imply exact energy conservation.

We test this distinction on the prior-informed complex fit from training
seed $17$ (GC17, defined in Section~\ref{sec:core-learning-comparisons}),
retaining its complete cosine and quartic terms. From initial state
$(0.1,0.06,0,0.02)$ we integrate to $\tau=10$ against a DOP853 reference
whose tolerance refinement changes each component by at most
$8.32\times10^{-14}$. Table~\ref{tab:core-integrator} confirms the
predicted orders.

\begin{table}[t]
\centering
\caption{Endpoint errors for the complete fitted GC17 energy. The state
error is Euclidean; orders use successive step halvings.}
\label{tab:core-integrator}
\setlength{\tabcolsep}{3.5pt}
\begin{tabular}{rcc}
\toprule
$\Delta$ & Midpoint error (order) & Central4 error (order)\\
\midrule
$0.08$ & $7.1351\times10^{-2}$ (---) & $1.8346\times10^{-3}$ (---)\\
$0.04$ & $1.9383\times10^{-2}$ ($1.8802$) & $1.1784\times10^{-4}$ ($3.9606$)\\
$0.02$ & $4.8966\times10^{-3}$ ($1.9849$) & $7.4149\times10^{-6}$ ($3.9902$)\\
$0.01$ & $1.2265\times10^{-3}$ ($1.9972$) & $4.6422\times10^{-7}$ ($3.9976$)\\
\bottomrule
\end{tabular}
\end{table}

The maximum implicit residual is $2.00\times10^{-15}$ and the maximum
symplectic defects are $3.34\times10^{-16}$ for Central4 and
$2.24\times10^{-16}$ for midpoint. An independent finite-difference
tangent agrees within $1.39\times10^{-12}$. These results confirm fourth-order convergence
and structure preservation at the measured solve accuracy. The comparison
uses equal step sizes; Central4 requires higher derivatives.
Supplementary Section~\ref{S-app:unified-integrator-check} gives the
solver tolerances and complete checks.

\subsection{Complex features and derivative approximation}
\label{sec:core-approximation}
Having controlled integration error, we next control the learned energy
and its derivatives. The complex form~\eqref{eq:core-complex-energy}
suggests retaining the magnetic cosine and learning a quartic residual.
For two modes, let $m_4(x)$ list the $35$ degree-four monomials,
$v_4$ contain the coefficients of $-(d^{\mathsf T}X)^4/24$, and the
$34$ columns of $U$ span its orthogonal complement. Define
\begin{align}
 F(x)&=\begin{pmatrix}g(d^{\mathsf T}X)/\|v_4\|_2\\
                           U^{\mathsf T}m_4(x)\end{pmatrix},\nonumber\\
 h_\eta(x)&=h_2(x)+\eta^{\mathsf T}F(x)
          =h_2+a\,g(d^{\mathsf T}X)+m_4(x)^{\mathsf T}Uc,
 \label{eq:unified-feature-readout}
\end{align}
where $\eta=(a\|v_4\|_2,c)$. This class contains the target at $a=1,c=0$
without duplicating its leading quartic direction. Substituting
\eqref{eq:core-complex-chart} realizes each monomial with complex affine
forms and their conjugates; square activations implement products via
$uv=[(u+v)^2-(u-v)^2]/4$. The cosine and these polynomial features give
a finite analytic complex network with a real linear readout and allow
momentum-dependent residuals.

More generally, increasing the polynomial degree approximates an
analytic target and each fixed order of its derivatives on a compact
polydisk inside the power-series convergence domain. The fixed class
above instead admits a quantitative coefficient bound. This derivative
requirement differs from the $L^2$ comparison for the particular
function classes in \citet[Theorem~2]{zhang}.

For an in-class target $h_*=h_2+\eta_*^{\mathsf T}F$, define
$B_r(D)=\max_{|\beta|\le r}\sup_{x\in D}\|\partial^\beta F(x)\|_2$.
For conservative derivative labels, let $\ell$ be the mean over $N$
observations of the squared Euclidean residual of each four-component
vector, and $\epsilon_{\rm rms}$ the corresponding label-error bound.
With $DF\in\mathbb R^{35\times4}$ and
$G=N^{-1}\sum_i DF(x_i)DF(x_i)^{\mathsf T}$, the triangle inequality gives
\begin{align}
 \|h_\eta-h_*\|_{C^r(D)}
 &\le B_r(D)\|\eta-\eta_*\|_2\nonumber\\
 &\le \frac{B_r(D)(\sqrt\ell+\epsilon_{\rm rms})}
              {\sqrt{\lambda_{\min}(G)}},
 \label{eq:core-feature-error}
\end{align}
provided $G$ is positive definite; the $C^r$ norm is the maximum over
scalar coordinate derivatives. Indeed, orthogonality of $J$ equates the
vector-field coefficient error to the derivative-feature error, and
the least eigenvalue of $G$ bounds its empirical squared norm from
below. This model-specific estimate includes nonzero fitting and label
errors. A target outside the class requires an additional approximation
defect in both the label budget and function error.
Section~\ref{sec:kam} uses these quantities to test KAM conditions.

\subsection{Matched identification and long-time prediction}
\label{sec:core-learning-comparisons}
Five paired comparisons separate representation from the mechanical
prior. R and C denote real and complex features; G fixes the physical
quadratic energy and B fits its $10$ coefficients. Thus GR and GC fit
$35$ real parameters, and BR and BC fit $45$, using the fixed
chart~\eqref{eq:unified-scale}. Within each prior choice, both
representations use identical function classes, data and regularizers
in common real coordinates.

For each seed in $\{17,43,91\}$, the two-mode Galerkin model generates
$12$ conservative $4$ s trajectories with $129$ observations each.
States are noiseless; vector-field components receive uniform noise
in $[-10^{-6},10^{-6}]$. Regularized least squares fits the readout,
and eight independently initialized $16$ s trajectories are held out
in full. Supplementary Section~\ref{S-sec:unified-learning-protocol}
specifies initial-state sampling, regularization and solver tolerances.

Abbreviating BR and BC as R and C, R$-$GR and C$-$GC measure the prior
effect, R$-$C and GR$-$GC compare representations, and
(GR$-$GC)$-$(R$-$C) measures their interaction. Differences are paired
within each training seed before computing means and sample standard
deviations.

\begin{table}[t]
\centering
\caption{Five paired contrasts of tip RMSE, using the same eight held-out
trajectories for each of three training seeds. Positive prior contrasts
favour G; positive representation contrasts favour C.}
\label{tab:core-learning}
\setlength{\tabcolsep}{3.5pt}
\begin{tabular}{lc}
\toprule
Contrast & Mean $\pm$ sample SD (mm)\\
\midrule
$\mathrm R-\mathrm{GR}$ & $(1.08\!\pm\!1.44)10^{-5}$\\
$\mathrm C-\mathrm{GC}$ & $(1.08\!\pm\!1.44)10^{-5}$\\
$\mathrm R-\mathrm C$ & $(-4.24\!\pm\!4.84)10^{-13}$\\
$\mathrm{GR}-\mathrm{GC}$ & $(0.75\!\pm\!1.99)10^{-15}$\\
$(\mathrm{GR}-\mathrm{GC})-(\mathrm R-\mathrm C)$ & $(4.25\!\pm\!4.82)10^{-13}$\\
\bottomrule
\end{tabular}
\end{table}

The mean tip RMSE is $2.28\times10^{-5}\,\mathrm{mm}$ for BR/BC and
$1.20\times10^{-5}\,\mathrm{mm}$ for GR/GC. The prior improves two seeds
and worsens the third. Representation and interaction contrasts are
smaller than the reference-solver refinement change,
$1.92\times10^{-8}\,\mathrm{mm}$; the matched real and complex fits
are numerically equivalent here. The three-seed summaries in
Table~\ref{tab:core-learning} are descriptive.
GC estimates $a=1.00326\pm0.00154$ with nonzero training residuals.
Supplementary Table~\ref{S-tab:unified-learning} gives slope and physical
target-energy errors as well as tip errors. The small errors reflect
an in-class target and accurately controlled synthetic data.

A separate $200$ s continuation spans $363$ first-mode periods and gives
GC tip RMSEs of $1.38\times10^{-4}$--$4.03\times10^{-4}\,\mathrm{mm}$,
maximum modal phase error $4.19\times10^{-4}\,\mathrm{rad}$, and
own-energy relative drift below $5.11\times10^{-9}$.
These DOP853 trajectories test long-time prediction, independently of
the Central4 test (supplementary Figure~\ref{S-fig:unified-long}).
Their finite-time agreement leads to the separate invariant-torus
question in Section~\ref{sec:kam}.

\subsection{Dissipation and independent finite-element comparisons}
\label{sec:core-physical-comparisons}
Driven tests first isolate the role of dissipation. Fitting the Rayleigh
coefficient on unforced trajectories gives tip RMSEs of
$6.40\times10^{-5}$--$1.37\times10^{-4}\,\mathrm{mm}$ on six held-out
forced trajectories; omitting it gives approximately $3.01\,\mathrm{mm}$.
The largest work-balance residual with dissipation is
$4.41\times10^{-17}\,\mathrm J$. These free-motion tests precede the
clearance constraints introduced in Section~\ref{sec:application}.

Independent COMSOL calculations then quantify reduction error.
On an $8$ s driven test, exact-cosine predictors with $2$, $4$ and $8$
modes give tip RMSEs of $0.0156168$, $0.00407422$ and
$0.00102484\,\mathrm{mm}$, respectively, versus
$0.00005614\,\mathrm{mm}$ for all $64$ modes.
The learned two-mode GC error, $0.0156171$--$0.0156177\,\mathrm{mm}$,
shows that reduction dominates this test. Supplementary
Section~\ref{S-sec:unified-comsol} gives the assembly and refinement checks.

Fitting position--velocity increments from separate full-FE COMSOL
trajectories gives $a=1.0223248822$ and test errors $0.14836$ and
$0.05140\,\mu\mathrm m$, versus $0.24172$ and $0.27459\,\mu\mathrm m$
for the exact two-mode predictor. The learned residual partly absorbs
reduction error (Supplementary Section~\ref{S-sec:comsol-identification}).
The KAM proof instead uses separate two-mode COMSOL observations,
so that its identification error can be bounded directly.

Unknown-parameter and model-mismatch tests further assess the feature
choice (Supplementary Section~\ref{S-sec:robustness-newphysics}).
For a soft rod with positive field gradient, adding omitted physical
features reduces tip RMSE from $182.988$ to $9.829\,\mu\mathrm m$;
a freely fitted quadratic model gives $7.258\,\mu\mathrm m$.
At $0.5\%$ state noise, the corresponding errors are $183.028$,
$9.757$ and $7.216\,\mu\mathrm m$. Thus a physical prior can limit
absorption of reduction error. Finally, on six one-second forecasts
along a slowly drifting full-FE trajectory, causal adaptation from past
states reduces mean error from $387.483$ to $9.867\,\mu\mathrm m$.
These nonautonomous forecasts have a separate empirical scope from the
stationary KAM result.
